\documentclass[12pt]{article}
%^ Start of document
\usepackage{times}  %Makes text TimesNewRoman
\usepackage{setspace} %Allows you to set single or double space 
\usepackage{biblatex} %Imports biblatex package
\usepackage{graphicx} % Required for inserting images
\usepackage{authblk}
\usepackage[colorlinks=true, urlcolor=blue, linkcolor=blue]{hyperref}
\usepackage{subfig}
\usepackage{float}
\usepackage{tabularx}
\usepackage{multirow}
\usepackage[paperheight=11in,
   paperwidth=8.5in,
   top=20mm,
   bottom=20mm,
   left=40mm,
   right=40mm]{geometry}
\usepackage{moresize}
\usepackage{tikz}
\usepackage{xcolor}
\usepackage{physics}
\usepackage{amssymb}
\usepackage{mathrsfs}
\usepackage{amsmath}
\usepackage[document]{ragged2e}
\usepackage{comment}
\usepackage{xcolor}
\addbibresource{mm.bib}
\begin{document}
\singlespacing  %Sets single spacing for whole document
\title{Magnetic Moment}

\author[1]{Zachary Tullier}
\affil[1]{Student\\Physics Department\\ University of Houston - Clear Lake \\Houston, Texas\\U.S}
\date{}
\maketitle

\section{Introduction}
    Atomic magnetic moments arise from the motion and spin of electrons. In an atom, each electron contributes an orbital magnetic moment (from its orbital angular momentum) and a spin magnetic moment (from its intrinsic spin). Quantum mechanically, the orbital magnetic moment operator is $\hat{\mu}_L = -\frac{e}{2m}\hat{L}$, so its $z$-component has eigenvalues $- \frac{e\hbar}{2m} m_l \equiv -\mu_B m_l$ (defining $\mu_B = \frac{e\hbar}{2m}$ as the Bohr magneton) \cite{PerturbationTheory}. Similarly, the electron’s spin contributes a moment $\hat{\mu}_S = -\frac{g_s e}{2m}\hat{S}$, where $g_s\approx 2$ is the electron $g$-factor for spin \cite{PerturbationTheory}. The total magnetic moment of an atom is the vector sum of all its orbital and spin moments. (Nuclei also have magnetic moments, but these are much smaller and often neglected in atomic magnetism.)
    
    \subsection{Non-Degenerate Perturbation Theory and Magnetic Moments}
        Perturbation theory deals with how a system’s properties change under a small additional interaction (a perturbation). In non-degenerate perturbation theory (no two states share the same unperturbed energy), the first-order energy correction to a state is given by the expectation value of the perturbing Hamiltonian in that state. For an atomic state with no degeneracy, one can directly compute shifts due to a perturbation like a magnetic field. For example, if we apply a weak static magnetic field $\mathbf{B}$, the perturbation Hamiltonian is $H' = -\boldsymbol{\mu}\cdot \mathbf{B}$ \cite{ZeemanW}. In a non-degenerate state $|\psi^{(0)}\rangle$, first-order perturbation theory gives an energy shift $\Delta E^{(1)} = \langle \psi^{(0)}| -\boldsymbol{\mu}\cdot \mathbf{B} |\psi^{(0)}\rangle$. This is essentially the Zeeman effect: the atomic energy levels shift in proportion to the component of the magnetic moment along $\mathbf{B}$. In fact, for an electron orbital with magnetic quantum number $m_l$, one finds $\Delta E = \mu_B m_l B$ to first order \cite{Zeeman}. Including spin, the shift generalizes to $\Delta E = \mu_B g_J m_J B$, where $m_J$ is the total angular momentum projection and $g_J$ is the Landé g-factor accounting for both spin and orbital contributions \cite{ZeemanW}. Thus, non-degenerate perturbation theory predicts a linear splitting of energy levels according to their magnetic moments when a field is applied. If the unperturbed state has a non-zero magnetic moment, it will shift in energy at first order; if it has no net moment (e.g. a singlet state with $L=0$, $S=0$), then the first-order shift can be zero and higher-order effects must be considered. Beyond energies, the wavefunction of the state is also perturbed. The first-order correction $|\psi^{(1)}\rangle$ is a sum over other (unperturbed) states, which means the perturbed state acquires a small mixture of other configurations. This can slightly alter expectation values of observables like the magnetic moment. However, for non-degenerate cases, these changes appear at least to second order in the perturbation. In summary, with non-degenerate perturbation theory we can reliably compute how a small magnetic perturbation (such as a weak field or a subtle spin interaction) shifts atomic energy levels and slightly modifies the state’s magnetic moment. This approach is valid as long as the perturbation does not connect to any degenerate levels at the same energy \cite{PerturbationTheory}.
        
    \subsection{Degenerate Perturbation Theory and Level Splitting}
        When the unperturbed atomic system has degenerate energy levels (multiple distinct states with the same energy), ordinary perturbation theory needs refinement. In degenerate perturbation theory, one must diagonalize the perturbation within the subspace of those degenerate states \cite{PerturbationTheory}. Physically, a perturbation will typically lift the degeneracy by splitting the energy levels and creating new linear combinations as eigenstates. The magnetic moment of the atom can be significantly affected by this process, since the perturbation chooses preferred orientations or couplings in the degenerate manifold. A classic example is the spin-orbit coupling inside atoms. Without spin-orbit interaction, states of the same principal quantum number $n$ with different $L$ (orbital) and $S$ (spin) orientations can be degenerate. The spin-orbit interaction $H_{S-O} = f(r), \mathbf{L}\cdot \mathbf{S}$ introduces a coupling between the electron’s orbital and spin magnetic moments \cite{PerturbationTheory}. This term mixes the degenerate $|L, m_L; S, m_S\rangle$ states. By switching to a basis of total angular momentum $J=L+S$ (i.e. coupling $L$ and $S$), one finds that $H_{S-O}$ becomes diagonal – this is effectively applying degenerate perturbation theory to find the “good” quantum states. The result is that each originally degenerate level splits into two (for $S=\tfrac{1}{2}$): one state with total $J = L+1/2$ and another with $J = L-1/2$. These have different energies given by the expectation of $\mathbf{L}\cdot \mathbf{S}$ in those states. In physical terms, the energy depends on the relative orientation of the orbital and spin magnetic moments (aligned vs. anti-aligned). One of the $J$ levels is shifted down in energy and the other up, removing the degeneracy and producing the fine structure of atomic spectra. The size of the splitting can be computed by first-order perturbation theory within the degenerate manifold, using the matrix elements of $\mathbf{L}\cdot \mathbf{S}$ \cite{PerturbationTheory}. Importantly, the new eigenstates (with definite $J$) have well-defined total magnetic moments and $g$-factors, different from the uncoupled ones, because the spin and orbital contributions are now entangled. Another example of degenerate perturbations is the Zeeman effect on a level with magnetic sublevels. In zero field, an atomic state with total angular momentum $J$ has $(2J+1)$ degenerate substates labeled by $m_J$. Turning on a magnetic field $B$ along (say) the $z$-axis adds $H_Z = -\mu\cdot B = \frac{eB}{2m}(L_z + 2S_z)$ \cite{ZeemanUT} as a perturbation. These $m_J$ states were degenerate, but $H_Z$ is diagonal in the $|J, m_J\rangle$ basis and directly splits them in energy. We can treat this by (degenerate) perturbation theory: the first-order shifts are $\Delta E = \langle J,m_J|H_Z|J,m_J\rangle$. This yields $\Delta E = \mu_B g_J B m_J$ as mentioned, so the $(2J+1)$ degenerate levels fan out into evenly spaced levels proportional to $m_J$ \cite{Zeeman}. Here the perturbation was already diagonal in the obvious degenerate basis, so the solution is straightforward (each $m$-state shifts independently without mixing). In more complicated cases (e.g. if $B$ is not aligned with a principal axis, or if there are accidental degeneracies between different $J$ levels), one would set up the Hamiltonian matrix in the degenerate subspace and diagonalize it to find the new eigenstates and energies. In all cases, degenerate perturbation theory ensures that the lifting of degeneracy is handled correctly, giving the new magnetic moment values for each split level.

    \subsection{First-Order vs. Second-Order Perturbative Effects on Magnetism}
        In perturbation theory, effects can appear at different orders in the perturbation strength (often denoted by a power of a small parameter). First-order effects are typically linear in the perturbation, while second-order effects are proportional to the perturbation squared (or involve two perturbative interactions). Both can influence atomic magnetic properties in different ways:
        \begin{enumerate}
            \item First-Order Perturbations: These give the dominant, linear response. If an atomic state has a permanent magnetic moment (e.g. due to unpaired spin or orbital angular momentum), a first-order perturbation will usually produce a direct change in energy or orientation. For instance, the Zeeman splitting discussed above is a first-order effect: energy levels shift linearly with $B$ according to $\Delta E^{(1)} \propto B$ \cite{Zeeman}. Likewise, the spin-orbit coupling term $H_{S-O} = f(r),\mathbf{L}\cdot\mathbf{S}$ produces a first-order fine-structure splitting of atomic levels proportional to the strength of that interaction \cite{PerturbationTheory}.First-order perturbations can also change the wavefunction at order $\lambda^1$, which in turn can lead to first-order changes in observable expectation values if the unperturbed expectation was non-zero. In general, first-order effects account for things like the immediate alignment of magnetic dipoles with a field (paramagnetism) and the primary shifts in spectra.
            
            \item Second-Order Perturbations: These are typically smaller corrections, but they can become important in cases where first-order effects cancel or are forbidden. Second-order energy corrections involve virtual transitions to intermediate states and are proportional to (perturbation strength)$^2$. In magnetism, second-order perturbations can induce a magnetic moment even when the atom has none initially, or they can modify the first-order results. A good example is the diamagnetic response of closed-shell atoms: if an atom has $J=0$ (total angular momentum zero) in its ground state, the first-order Zeeman shift is zero (since $m_J=0$ gives $\Delta E^{(1)}=0$). However, a magnetic field will still induce currents in the electron cloud. This appears as a second-order effect: the field mixes the ground state with excited states, leading to a small induced magnetic moment opposite to the field (Lenz’s law) and an energy shift proportional to $B^2$ (the quadratic Zeeman effect). For example, in the hydrogen 1s state (which has no orbital angular momentum), the leading Zeeman effect comes from a second-order term resulting in $\Delta E \propto +B^2$ (a small diamagnetic shift) \cite{QZeeman}. Second-order perturbation theory also explains Van Vleck paramagnetism, a subtle form of magnetism in certain ions: there, a $J=0$ ground state can have a weak, temperature-independent paramagnetism due to excited-state mixing in second order \cite{Magnetism}. Generally, second-order perturbations refine the magnetic properties by accounting for induced moments and polarization effects, and they become crucial when first-order terms vanish or when higher precision is needed. They also give corrections like the quadratic Zeeman splitting (important at high fields or in precision spectroscopy) and small shifts in $g$-factors due to state mixing.
        \end{enumerate}
 
        In summary, first-order perturbations tend to dominate the observable magnetic effects (linear energy shifts, level splittings, etc.), while second-order perturbations provide smaller corrections (often quadratic in field) that can induce magnetism in otherwise non-magnetic states or produce fine adjustments to energy levels.

        \subsection{Zeeman Effect and Magnetic Perturbations}
        
            The Zeeman effect is the prototypical example of an external perturbation on an atomic magnetic moment. Here, an external magnetic field $\mathbf{B}$ is applied to an atom, and the interaction $H_Z = -\boldsymbol{\mu}\cdot \mathbf{B}$ perturbs the atomic Hamiltonian. In vector form, the atomic magnetic moment can be written as $\boldsymbol{\mu} \approx -\frac{\mu_B g \mathbf{J}}{\hbar}$ (for the electronic contribution, neglecting nuclear magnetism) \cite{ZeemanW}. Thus, $H_Z = -\boldsymbol{\mu}\cdot \mathbf{B} = \mu_B g \frac{\mathbf{J}\cdot \mathbf{B}}{\hbar}$. If $B$ is along $z$, this becomes $H_Z = \mu_B g J_z B/\hbar$, which is diagonal in $m_J$. First-order perturbation theory then immediately gives the energy shift $\Delta E = \mu_B g m_J B$ \cite{ZeemanW} for a state labeled by $J, m_J$. Empirically, this explains the splitting of spectral lines into multiple components in a magnetic field – an observation made by Pieter Zeeman in 1896 \cite{ZeemanW}. For systems where $g$ is known (Landé g-factor for a given level), the Zeeman splitting is linear in $B$ (the linear Zeeman effect or first-order Zeeman effect). This is often sufficient for weak fields. However, for very strong fields or certain special cases, second-order (quadratic) Zeeman effects appear. One scenario is when two levels become quasi-degenerate due to a field – then second-order mixing can shift them. Another scenario is when the linear term is zero (as discussed, $J=0$ states have no first-order Zeeman splitting). In such cases, higher-order perturbation must be included, leading to energy shifts like $\Delta E^{(2)} \propto B^2$. This quadratic Zeeman effect is usually much smaller than the linear term and becomes noticeable only at high fields (or through very precise measurements) \cite{QZeemanE}. For example, in hydrogen, the $n=1$ ground state’s only Zeeman shift comes from a second-order effect involving virtual transitions to $n=2$ states, yielding a tiny quadratic dependence on $B$. Aside from shifting energies, the Zeeman perturbation also influences the orientation of atomic magnetic moments. In weak fields, the atomic moments precess around the field (Larmor precession) but remain largely in their eigenstates aligned with $J$. In stronger fields, if the Zeeman splitting exceeds other splittings (like fine structure), the coupling of $L$ and $S$ can be broken (the Paschen-Back regime), which is essentially a non-perturbative regime beyond the scope of first-order theory. Perturbatively, one can approach that regime by summing higher orders or using degenerate theory carefully, but typically one switches basis to $L_z, S_z$ as the good quantum numbers when $B$ is very large. This highlights that perturbation theory has limits: it works best when the perturbing field is “small” compared to the intrinsic atomic fields (such as the Coulomb field establishing the level spacings or the spin-orbit interaction). For normal laboratory fields, the first-order Zeeman formula is usually accurate, and second-order corrections are tiny. Important to note: The Zeeman effect demonstrates how an external perturbation can directly measure and manipulate atomic magnetic moments. By analyzing Zeeman splittings, one can determine $g$-factors and infer the composition of an atom’s magnetic moment (spin vs orbital) \cite{ZeemanUT}. In modern research, this is fundamental for techniques like laser spectroscopy of atoms, magnetic tuning of atomic energy levels (e.g. in quantum computing with trapped atoms/ions), and probing astrophysical magnetic fields via spectral line splitting. All of these applications rely on the basic perturbative understanding of how atomic magnetic moments respond to external fields.

\printbibliography    
\end{document}
